Nearly $30$ years after the introduction of the first lattice-based cryptosystem~\cite{Ajtai96}, lattice-based cryptography is now in widespread deployment. A major reason for this is the threat of large-scale quantum computers and the harvest now, decrypt later problem. Following NIST standardization of the Kyber key exchange mechanism and Falcon signature, we have seen them integrated in applications used by billions daily such as TLS, SSH, Signal and iMessage, Zoom, Wireguard and IKEv2, and MLS. We have also seen increasing real-world deployment of fully homomorphic encryption, all of which is lattice-based.

%For many areas of post-quantum cryptography, the most efficient constructions have come from lattices. One area where this is not the case is succinct arguments, particularly with regards to verifier time. For example, for circuits of size $2^{24}$ state-of-the-art hash-based arguments such as WHIR report verification times of $0.6$ms, while lattice-based systems such as Greyhound~\cite{C:NguSei24} take $130$ms--an important difference for verifier constrained environments like blockchains and embedded devices far less powerful than the 64-core CPU used for the benchmarks above. Recent lattice-based arguments improve on Greyhound's $O_{\lambda}(\sqrt{n})$ verifier, but unfortunately lack implementations. This includes the polynomial commitment (PC) of Cini, Malavolta, Nguyen, and Wee~\cite{C:CMNW24} and Hyperwolf~\cite{Hyperwolf}, both of which boast $O_{\lambda}(\log n)$ verifiers. It is plausible that the concrete verifier time of these schemes will continue to lag behind hash-based competitors, as they both have at least a quartic dependence on the security parameter~\pnote{confirm this}, while for hash-based schemes this dependence is strictly linear.

As lattice-based cryptography and succinct arguments have both realized increased use, naturally the interest in proving statements about lattice-based cryptography has grown as well. One important use case is building efficient versions of more complex primitives by verifying simpler ones using succinct non-interactive arguments of knowledge (SNARKs). In the lattice setting this approach has been particularly successful for constructing various types of signatures--aggregate, blind, group, multi, threshold, and ring--from standard signatures~\cite{PQCRYPTO:BCOS20,AC:LyuNgu22,C:AABKT24}. SNARKs have also been used to construct building blocks needed for other schemes such as...

Folklore suggests that lattice-based arguments are the natural candidate for verifying lattice-based cryptography. This work studies using hash-based arguments instead. One of our contributions is the first construction of a hash-based PC for proving evaluations directly over the cyclotomic rings used by all modern lattice-based schemes. This approach does not actually suggest a contradiction, as the crux of the folklore argument was that lattice-based arguments natively proved statements over cyclotomic rings. We've simply shown how to construct linear codes over cyclotomic rings with the necessary proximity gaps required for building succinct PCs.

One compelling counterargument to the folklore is that arguments that work directly over cyclotomic rings $Z_{q}[X]/(\Phi(X))$ pay a linear proof size cost in the ring dimension $n$, where concretely we often have $n=512$ or larger. On the other hand, prior work simulating ring arithmetic in $\field_q$ has incurred a quadratic prover complexity depdendence on $n$~\cite{PQCRYPTO:BCOS20}. An innovative new result of Huang, Mao, and Zhang~\cite{EPRINT:HuaMaoZha25} shows it's possible to avoid both costs via a new \emph{ring-switching} argument, which recasts the statement over a Galois ring for which they present a new hash-based PC. Our work provides another way past this dichotomy by introducing an optimal $O(n \log n)$ prover complexity interactive oracle proof (IOP) for simulating ring multiplication over fields, which works with any field-agnostic multilinear PC. This avoids the constraint blowup of the ring-switching approach...

One obstacle we encounter when simulating ring arithmetic over $\field_q$ is that the small moduli used in many lattice-based constructions do not confer sufficient soundness when used with IOPs and hash-based PCs over the field. Numerous recent works have introduced methods for working over small fields~\cite{ITCS:BGKS20,EPRINT:StarkWare21,EPRINT:HabLevPap24,EC:DiaPos25,EPRINT:DiaPos24,EPRINT:BagDomTha24}, all of which use some type of embedding of the witness in a sufficiently large extension field. Diamond and Posen present a particularly elegant way work over binary extension fields with essentially zero embedding overhead~\cite{EPRINT:BagDomTha24}. We extend their aproach to work over fields of odd characteristic.

\pnote{Implementation details go here.}


%%%%%%%%%%%%%%%%%%%%%%%%%%%%%%%%%%%%%%%%%%%%%%%%%%%%%%%%%%%%%%%%%%%%%%%%%%%%%%%%
%%%%%%%%%%%%%%%%%%%%%%%%%%%%%%%%%%%%%%%%%%%%%%%%%%%%%%%%%%%%%%%%%%%%%%%%%%%%%%%%
\subsection{Contributions}
\label{sec:contributions}

This work then makes the following contributions:
\begin{enumerate}[nosep]
  \item We show how any linear $[n,k,d]$-code over a field $\field_q$ gives a linear $[n,k,d]$-code for all cyclotomic rings of the form $\integers_q[X]/(\Phi_{m}(X))$. We then further prove that linear codes over all finite rings $\ring$ (and even certain infinite rings) satisfy a $(d/3,E)$-proximity gap in the unique decoding radius regime, where $E$ is equal to the GCD of the order of its residue fields (for cyclotomic rings $E$ is the order of its largest exceptional set minus one). This result enables us to construct an IOP of proximity to the ring, which we in turn use to build a hash-based multilinear PC over cyclotomic rings.
  \item We present a new sumcheck-based IOP for $f$-circulant matrix multiplication. The $f=-1$ case corresponds to multiplication in power-of-two cyclotomic rings. If $q$ is chosen to support a NTT over the ring as is generally the case, we . We also present a slightly less efficient IOP for multiplication in odd cyclotomics using the Chinese remainder transform.
  \item We generalize the packing technique of Diamond and Posen~\cite{EPRINT:BagDomTha24} to fields of odd characteristic.
  \item \pnote{Implementation details go here.}
\end{enumerate}
